
% \subsection{Ánh xạ tuyến tính}
% \noindent Giả sử \(U\) và \(W\) là các không gian véc tơ trên trường \(\K\).
% \begin{define}[\textit{Ánh xạ tuyến tính}]
%     Ánh xạ \(f: U \to W\) được gọi là một \textit{ánh xạ tuyến tính}, nếu
%     \begin{align*}
%         f(u + v) &= f(u) + f(v) \\ 
%         f(\alpha u) &= \alpha f(u)
%     \end{align*}
%     với mọi \(u, v \in V\) và mọi vô hướng \(\alpha \in \K\).
%     \label{def:linear_map} 
% \end{define}
% Hai điều kiện trong định nghĩa ánh xạ tuyến tính là tương đương với điều kiện sau:
% \begin{align*}
%     f(\alpha u + \beta v) = \alpha f(u) + \beta f(v)
% \end{align*}
% Một số tính chất của ánh xạ tuyến tính được suy ra từ định nghĩa. Giả sử \(f:V \to W\) là một ánh xạ tuyến tính, khi đó ta có:
% \begin{enumerate}[label=(\roman*)]
%     \item \(f(0) = 0\).

%     Thật vậy, \(f(0) = f(0+0) = f(0) + f(0)\). Theo luật giản ước, \(f(0)=0\).
%     \item \(f(-u) = -f(u)\), \(\forall u \in V\).
    
%     Thật vậy, \(f(u)+f(-u) = f(u+(-u)) = f(0) = 0\). Do đó, \(f(-u)=-f(u)\).

%     \item \( f\left(\displaystyle{\sum_{i=1}^n \alpha_i u_i}\right) = \displaystyle{\sum_{i=1}^{n} \alpha_i f(u_i)} \), \(\forall u_i \in V\), \(\alpha_i \in \K\).
    
%     Từ định nghĩa \eqref{def:linear_map}, ta suy ra công thức trên là đúng.
%     Giả sử công thức đúng với \(n=k\), tức là:

%     \( f\left(\displaystyle{\sum_{i=1}^k \alpha_i u_i}\right) = \displaystyle{\sum_{i=1}^{k} \alpha_i f(u_i)} \), \(\forall u_i \in V\), \(\alpha_i \in \K\).

%     Ta cần chứng minh:

%     \( f\left(\displaystyle{\sum_{i=1}^{k+1} \alpha_i u_i}\right) = \displaystyle{\sum_{i=1}^{k+1} \alpha_i f(u_i)} \), \(\forall u_i \in V\), \(\alpha_i \in \K\).


%     Thật vậy, ta có:
%     \begin{align*}
%         f\left(\displaystyle{\sum_{i=1}^{k+1} \alpha_i u_i}\right) &=   f\left(\displaystyle{\sum_{i=1}^{k}\alpha_i u_i + \alpha_{k+1} u_{k+1}}\right)\\ &= f\left(\displaystyle{\sum_{i=1}^{k} \alpha_i u_i}\right) + f(\alpha_{k+1} u_{k+1}) \\ &=\displaystyle{\sum_{i=1}^{k+1} \alpha_i f(u_i)}.
%     \end{align*}
% \end{enumerate}